\subsection*{一、$C$ 隨光學厚度下降的原因}
光學厚度較小時，多數光子只經歷一次散射，角度分布仍接近單一粒子的散射分布，前向的 $30^\circ$ 強度遠高於後向的 $125^\circ$ 強度，因此 $C$ 較大。隨著 $\tau$ 增加，光子經歷更多次散射並逐漸失去原有的前進方向，前向峰被壓平，$C$ 隨之下降。

方向記憶消失所需的特徵距離約為傳輸平均自由徑 $l^*=l/(1-g)$，其中 $l$ 為平均自由徑（Pine et al., 1990）。本研究中，$2.16\,\mu\mathrm{m}$ 微球的異向性參數 $g=0.907$，高於 $0.492\,\mu\mathrm{m}$ 微球的 $g=0.818$；前者單次散射的平均偏轉較小，因此在相同 $\tau$ 下，$C$ 下降較慢。

以 $\ln C$ 對 $\tau$ 作直線擬合，實驗斜率在 $2.16\,\mu\mathrm{m}$ 與 $0.492\,\mu\mathrm{m}$ 微球分別為 $-0.224\pm0.007$ 與 $-0.254\pm0.019$；模擬斜率則分別為 $-0.267\pm0.021$ 與 $-0.298\pm0.032$。兩者均支持較大粒徑的下降較慢，但差異僅約 $10\%$ 至 $15\%$。若下降速率只由 $1-g$ 決定，依兩者 $1-g$ 約 $1.95$ 倍的差異，較大粒徑的斜率絕對值應約為另一者的一半；實際結果並非如此。以 $\tau(1-g)$ 作為橫軸時，兩組曲線也未重合。這表示在本研究範圍內，平均異向性參數能說明趨勢方向，卻不足以單獨決定 $C$ 的下降速率；散射相位函數的其他角度細節仍有影響。此範圍的 $\tau(1-g)\leq1.5$，光子尚未進入完全擴散狀態。

\subsection*{二、實驗與模擬的差異}
當 $\tau\geq2$ 時，兩種粒徑的實驗 $C$ 均高於模擬；其中 $0.492\,\mu\mathrm{m}$ 在 $\tau=6$、$8$ 時的實驗／模擬比分別為 $1.31$、$1.29$。樣品管中的背景光應介於完全不扣空白與扣除全部空白兩種極端情形之間；檢查這兩種扣除方式後，差異仍無法消除。可能原因是裝置反光：比色槽周圍的樂高零件與隨旋轉臂移動的線材，被樣品的散射光照亮後再反射進感測器（見參、十一）。圖~\ref{fig:shape} 顯示 $\tau\geq4$ 時偏差集中在前向，與此機制的方向一致。純水空白在臂角約 $40^\circ$ 有一個窄峰，對應線材橫過出射光束的位置，表示線材確實會被照到；但此峰僅約 $700$ 至 $1000$ 計數，遠小於 $2.16\,\mu\mathrm{m}$ 在 $\tau=4$、$6$ 時前向多出的約 $2500$ 至 $5900$ 計數，因此線材應只占一部分，其餘可能來自距比色槽更近的樂高零件。兩者的貢獻需以分步遮黑逐項量測：先包黑線材，再於管座與底座貼上黑絨布。

\subsection*{三、單一兩角強度比的限制}
雖然相同 $\tau$ 下較大粒徑的 $C$ 較高，兩種粒徑的 $C$ 數值範圍仍互有重疊。例如，$0.492\,\mu\mathrm{m}$ 在 $\tau=0.5$ 時的 $C$ 約為 $12.1$，也可能對應到較大粒徑在較高光學厚度時的數值。原稿列出的 12 個樣品中有 9 個，存在不同粒徑與不同 $\tau$ 得到相同 $C$ 的情形。因此，單一 $C$ 無法同時唯一決定粒徑與濃度；若要由 $C$ 推估濃度，需先知道粒徑，或加入其他獨立量測資訊。

\subsection*{四、未來工作}
\begin{enumerate}
    \item 每個濃度配製多管並重複量測，以評估配製誤差及樣品放置時間的影響。
    \item 量測雷射偏振方向與光束位置，並在模擬中納入液面及比色槽底部，釐清 $\tau\geq2$ 時實驗值高於模擬值的原因。
    \item 建立 $C$ 隨 $\tau$ 變化的經驗關係式，在粒徑已知時由 $C$ 反推光學厚度與濃度。
    \item 結合兩組以上角度比，或加入單一角度強度，評估粒徑未知時同時估算粒徑與濃度的可行性。
    \item 探討能否利用 $C$ 判斷單一角度強度位於濃度響應曲線峰值的哪一側，改善單角度讀值的濃度判讀。
    \item 擴大樣品範圍，加入更多粒徑、聚乙烯（PE）、聚對苯二甲酸乙二酯（PET）、混合粒徑及不規則形狀微塑膠，檢驗方法適用範圍及 $C$ 下降速率與 $g$ 的關係。
    \item 朝低成本、可於現場快速估算水中微塑膠粒徑與濃度的量測裝置發展。
\end{enumerate}
